% RAKL_V3_MANUSCRIPT_UPDATE_20260811
\section{Directional witnesses as RAKL v3 experience-transfer relations}
\label{sec:v3-transfer}

RAKL v3 makes experience persistence an implemented architectural fact but leaves transfer validity as an evidence question. A source \texttt{TaskEpisode} says what happened in one registered context. A versioned \texttt{Lesson} is a scoped abstraction over one or more episodes. Neither object is automatically applicable to a new target. The role of the directional structural witness is to supply an explicit candidate relation
\[
W_{s\rightarrow t}^{(q)}:
(E_s,L_s,\gamma_s)
\rightsquigarrow
(P_t,q,\gamma_t)
\]
whose obligations state which roles and invariants are preserved, which properties are not preserved, and which target boundaries would invalidate reuse.

This makes the witness directional in two senses. First, the target quantity of interest determines which source structure is load-bearing. Second, a valid mapping from source to target need not license the reverse mapping because scope, information loss or boundary conditions can be asymmetric. A witness is therefore not an embedding similarity score and not an equivalence class over tasks.

\subsection{Transfer can alter search priority without creating authority}
When a witness passes the registered structural and boundary gate, v3 may use the linked episode or lesson to increase retrieval priority, activate an already promoted tool, or place an operator path earlier in the search queue. The permitted inference is
\[
\text{witnessed applicability}
\Rightarrow
\text{eligible experience reuse},
\]
not
\[
\text{witnessed applicability}
\Rightarrow
\text{scientific claim true}.
\]
Scientific authority remains owned by the target problem's verification path. This is the cross-domain specialization of the framework invariant that experience-conditioned routing is not epistemic authority.

\subsection{Negative transfer is learned through an evidence ladder}
A failed transfer attempt first creates a non-success task episode. It does not immediately establish that the witness class, operator or source lesson is invalid. The system records the failure as observed-only, then permits competing diagnoses such as context mismatch, broken invariant, implementation defect, missing evidence or an unrelated downstream failure. Only discriminating evidence can promote a supported diagnosis, and only fresh replay/transfer can turn that diagnosis into a reusable boundary lesson.

The resulting discipline prevents one near-miss from becoming a global blacklist while still allowing repeated boundary-sensitive failures to improve later routing. It also makes negative-transfer history addressable: a future witness can explicitly inherit a known non-preserved property or target boundary rather than merely receiving a lower opaque score.

\subsection{Connection to RAKL-triviality}
V3 classifies a verified solution as \texttt{TRANSFER\_NOVEL} when no new primitive was invented but pre-existing problem-solving structure was transported into a new context by an explicit mapping witness. The directional witness studied here is a candidate evidence object for that classification. A solution should receive \texttt{TRANSFER\_NOVEL} only when the transfer relation is explicit, the target solution is independently verified, and resource ancestry shows that no new representation/operator/ontology was required.

This gives Paper III a narrower post-v3 claim. The framework already possesses a persistent experience substrate; this paper does not claim novelty for persistence itself. It asks whether one scientifically scoped, boundary-aware relation can make cross-domain reuse safer and more measurable than semantic resemblance alone. The natural-domain, strong-control and independent-annotation gates remain necessary before that question can receive an empirical answer.
